\begin{abstract}
Adapting a vision-language-action (VLA) model to a new robot can require learning a new action module. The model must learn how to control the new embodiment, often before successful task executions are available. We propose a simple training sequence for this setting. First, reinforcement learning with dense rewards teaches the individual skills of a manipulation task by providing feedback on partial progress. Next, binary rewards focus each skill on successful completion, while self-imitation consolidates successful experience. Finally, staged reinforcement learning combines these skills into one policy that executes the complete instruction. We study this approach by learning a residual action module over pretrained SmolVLA on a cable-driven manipulator, chosen as an example of an unusual robot embodiment. The adapted controller achieves 61.43\% strict full-task success in simulation development evaluation. A complementary OpenVLA-OFT experiment on LIBERO examines stage-aware feedback for an already trained manipulator. \assumed{Physical experiments demonstrate complete object manipulation on the target robot.} Together, the studies investigate how to build useful robot behavior from pretrained VLA representations through a sequence of learnable steps.
\end{abstract}
